\documentclass[10pt,a4paper]{amsart}
\usepackage[margin=19mm]{geometry}
\usepackage{amsmath,amssymb,mathtools}
\usepackage[scr=rsfs,cal=euler]{mathalfa}
\usepackage{xfrac}
\usepackage[unicode,hidelinks,bookmarks=false]{hyperref}
\newcommand{\ie}{i.e.\ }
\newcommand{\cf}{cf.\ }
\newcommand{\resp}{respectively\ }
\newcommand{\Subsection}{\subsection}
\providecommand{\nobreakdash}{\nobreak\hbox{-}}
\setlength{\emergencystretch}{2em}
\multlinegap0pt
\usepackage{bm}
\usepackage{bbm}
\usepackage{dsfont}
\usepackage{xspace}
\usepackage{cancel}
\usepackage{mathtools}
\usepackage{amsthm}
\makeatletter
\@ifundefined{lem}{\newtheorem{lem}{Lemma}}{}
\@ifundefined{prop}{\newtheorem{prop}[lem]{Proposition}}{}
\makeatother
\newcommand{\Q}{{\mathbb Q}}
\newcommand{\eps}{\varepsilon}
\newcommand{\gen}{{\operatorname{gen}}}
\newcommand{\rsp}{\raisebox{0em}[2.7ex][1.3ex]{\rule{0em}{2ex} }}
\title[Octic Hilbert 2-class fields of real quadratic fields with d]{Octic Hilbert 2-class fields of real quadratic fields with discriminant 8p}
\author{Franz Lemmermeyer}
\date{Source: arXiv:2510.10295v1 (CC BY 4.0)}
\begin{document}
\maketitle

\noindent\emph{Source and licence.} Octic Hilbert 2-class fields of real quadratic fields with discriminant 8p. Version arXiv:2510.10295v1 (CC BY 4.0), distributed under the Creative Commons Attribution 4.0 International licence, as recorded in the arXiv licence metadata for this exact version and shown on the abstract page https://arxiv.org/abs/2510.10295v1. Full rights notice: licenses/arxiv-2510.10295-CC-BY-4.0.txt.\par\smallskip

\noindent\textit{Editorial scope.} Counted unit: the Prop whose source label is \texttt{None} in the article (source0.tex lines 133--145, source label \texttt{None}), delivered together with the article's own complete original proof of it (source0.tex lines 147--169, one \texttt{proof} environment of 23 lines). No mathematics is written, repaired or re-derived: statement and proof are the article's own text line for line. Required context retained uncounted: Pr01 (lines 172--178); Lem1 (lines 219--223). Every definition, display and label of those lines is retained unchanged. Omitted from this package and not redistributed: 1--132, 146--146, 170--171, 179--218, 224--561. Nothing inside a retained excerpt is omitted or altered: every retained body is delivered exactly as the article's lines are written, so no omission receipt of any kind accompanies this package. Citations: the retained excerpts carry no citation command at all, so no bibliography entry of the article is transcribed and no reference is redistributed. Reference adaptation: none was needed and none was made; every reference inside the delivered unit resolves inside it, so no unresolved reference or citation is produced. Printed slips kept literal: no printed word, symbol, hypothesis or number of the counted statement or of its proof was altered. Layout and class: the article's own class is \texttt{amsart} with packages including amsmath, amssymb, amscd, amsthm, tikz, hyperref; the wrapper is the lane's pinned \texttt{amsart} editorial wrapper at 10pt on A4, which restates the theorem-like environments the retained text needs and adds the article's own macro definitions that the retained text needs (\texttt{\textbackslash Q}, \texttt{\textbackslash eps}, \texttt{\textbackslash gen}, \texttt{\textbackslash rsp}), with extra wrapper packages (bm, bbm, dsfont, xspace, cancel, mathtools). The delivered numbering is the unit's own. Assets: no retained span contains a figure environment, a graphics-inclusion command, a PDF-page inclusion, a table or a TikZ picture; no image file is copied into this package, the package declares no asset dependency and no image file is redistributed with it. Licence: version arXiv:2510.10295v1 of this article is distributed under the Creative Commons Attribution 4.0 International licence, as recorded in the arXiv licence metadata for this exact version and shown on the abstract page https://arxiv.org/abs/2510.10295v1. A full rights notice is delivered at licenses/arxiv-2510.10295-CC-BY-4.0.txt. Characters: every retained excerpt is pure ASCII, so the delivered engine, XeLaTeX with its default Latin Modern text font, reports no missing character.
\begin{prop}\label{Pr01}
  Let $p$ be an odd prime number. The class number $h^+$ in the strict
  sense of $\Q(\sqrt{2p}\,)$ is divisible by $8$ if and only if
  $p \equiv 1 \bmod 16$ and $(\frac 2p)_4 = +1$.
  We have $h \equiv 4 \bmod 8$ and $N\eps_{2p} = -1$  if and only if
  $p \equiv 9 \bmod 16$ and $(\frac 2p)_4 = -1$.
\end{prop}

\begin{lem}\label{Lem1}
  Let $p = e^2 - 2f^2 \equiv 1 \bmod 8$ be a prime and assume that
  $f \equiv 2 \bmod 4$.  Then
  $$ \Big(\frac p2 \Big)_4 = \Big(\frac p2 \Big)_4 = \Big(\frac{-2}e \Big). $$
\end{lem}

\begin{prop}
  Let $p \equiv 1 \bmod 8$ be a prime number. Then $p = e^2 - 2f^2$
  for integers $e$ and $f > 0$ with $e \equiv 3 \bmod 4$ and
  $f \equiv 2 \bmod 4$.

  $$ \begin{array}{lc|l}
    \rsp \text{class number} &  N\eps_{2p} & \text{conditions} \\ \hline
    \rsp  h \equiv 2 \bmod 4 & +1    & e < 0 \\
    \rsp  h \equiv 4 \bmod 8 & -1    & e > 0, \ e \equiv 7 \bmod 8 \\
    \rsp  h \equiv 4 \bmod 8 & +1    & e > 0, \ e \equiv 3 \bmod 8 \\
    \rsp  h \equiv 0 \bmod 8 & \pm 1 & e > 0, \ e \equiv 3 \bmod 8
  \end{array} $$
\end{prop}

\begin{proof}
  The extension $k(\sqrt{\alpha}\,)$ for $\alpha = e + f\sqrt{2}$
  is a cyclic quartic extension of $k = \Q(\sqrt{2p}\,)$ containing
  $k_\gen = \Q(\sqrt{2},\sqrt{2p}\,)$. Since
  $N(\alpha) = \alpha\alpha' = e^2 - 2f^2 = p$, it can only ramify above
  the primes dividing $2 p \infty$. Since
  $k(\sqrt{\alpha}\,) = k(\sqrt{\alpha'}\,)$ and since $\alpha$ and $\alpha'$
  are coprime,  it can only ramify at $2$ and $\infty$.

  Since $\alpha \equiv 3+2\sqrt{2} = (1+\sqrt{2}\,)^2 \bmod 4$, the
  extension $k(\sqrt{\alpha}\,)$ is a cyclic quartic extension
  unramified except possibly at the infinite primes.

  Thus if $e < 0$, then the $2$-class number of $k$ is $h = 2$ and the
  $2$-class number of $k$ in the strict sense is $h^+ = 4$; this proves the
  first line.

  If $e > 0$ and $e \equiv 7 \bmod 8$, then $(\frac p2)_4 = (\frac2p)_4 = -1$
  by Lemma~\ref{Lem1} below, hence $h = h^+ = 4$ by Prop.~\ref{Pr01}.

  If $e > 0$ and $e \equiv 3 \bmod 8$, then $(\frac p2)_4 = (\frac2p)_4 = +1$,
  hence $h^+$ is divisible by $8$. 
\end{proof}

\end{document}
